%
% Tex input file for CGNS SIDS-to-HDF File Mapping Manual
%
% To generate a DVI file, then a PostScript file named filemap_hdf.ps:
%
%     latex filemap_hdf.tex
%     dvips filemap_hdf.dvi -o
%
% To generate a PDF file named filemap_hdf.pdf:
%
%     pdflatex filemap_hdf.tex


\documentclass[twoside,fleqn]{article}

%
% Packages
%
\usepackage{indentfirst}		% indent first paragraph in sections
\usepackage[tbtags]{amsmath}		% use ams math package
\usepackage{graphicx}			% for including eps files
\usepackage{lscape}			% for landscape stuff
\usepackage[bf]{caption}		% for more flexible captions
   \setlength{\captionmargin}{18pt}	% caption margins
\usepackage{array}			% use new array and table package
   \setlength{\extrarowheight}{2pt}	% increase row heights in tables
\usepackage{longtable}			% use longtable package
\usepackage{alltt}			% verbatim input with font changes
\usepackage{natbib}			% for author-year citations
   \setlength{\bibhang}{0pt}		% don't use hanging indent format
\usepackage[normalem]{ulem}		% allow underlining
\usepackage{calc}			% calculation package
\usepackage{ifthen}			% control structures
\usepackage{fancyhdr}			% fancy headers and footers
   \fancyhf{}
   \renewcommand{\headrulewidth}{0pt}
   \renewcommand{\footrulewidth}{0pt}
\usepackage{titling}			% more title page control
\usepackage{xspace}			% common-sense spacing after text macro
\usepackage{mdwlist}			% more flexible description lists
\usepackage{epic}			% extended picture environment
%\usepackage[pdftex]{color}		% colors
\usepackage{color}			% colors
   \definecolor{subcolor}{rgb}{0.809,0.852,1.0}
\usepackage{pdfcolmk}			% for colors across page breaks
%\usepackage[T1]{fontenc}		% for better-looking "_" and tt "{"
\usepackage{ae}				% used instead of above for PDF output
\usepackage{hyperref}			% for hypertext links in PDF
   \hypersetup{letterpaper,plainpages=false,
               pdftitle={HDF5 Implementation},
               pdfauthor={CGNS Project Group},
               colorlinks,
               linkcolor=blue,citecolor=blue,filecolor=red,pagecolor=blue,
               urlcolor=red}
   \renewcommand{\sectionautorefname}{Section}
   \renewcommand{\subsectionautorefname}{Section}
   \renewcommand{\subsubsectionautorefname}{Section}
   \renewcommand{\paragraphautorefname}{Section}

%
% Page layout
%
\normalsize
\setlength{\oddsidemargin}{0.5in}	% left margins-1.0in for odd/even pages
\setlength{\evensidemargin}{0.0in}
\setlength{\textwidth}{6.0in}		% text width
\setlength{\topmargin}{-0.5in}		% top margin-1.0in
\setlength{\headheight}{0.25in}		% header space parameters
\setlength{\headsep}{0.5in}
\setlength{\textheight}{8.5in}		% text height
\setlength{\topskip}{\baselineskip}	% dist from top of body to 1st baseline
\setlength{\footskip}{0.75in}		% dist from bottom of body to footer
\raggedbottom				% to avoid stretching vertical space

\pagestyle{fancy}			% using fancyhdr package
\fancypagestyle{plain}{%		% page # in corner for plain style
   \fancyhf{}
   \fancyfoot[LE,RO]{\bfseries \thepage}}
\addtocounter{secnumdepth}{1}		% number \paragraph's
\addtocounter{tocdepth}{1}		%    and include them in TOC

%
% Misc spacing parameters
%
\setlength{\parskip}{0.5\baselineskip}	% blank space between paragraphs
\setlength{\doublerulesep}{0pt}		% for wider \hlines
\setlength{\fboxsep}{5pt}		% for more margin in boxes
\newlength{\saveparindent}		% save for use where LaTeX changes them
\newlength{\saveparskip}		% save for use where LaTeX changes them
\newlength{\savebaselineskip}
\setlength{\saveparindent}{\parindent}
\setlength{\saveparskip}{\parskip}
\setlength{\savebaselineskip}{\baselineskip}
\newlength{\tmplength}			% for temp use wherever needed
\newlength{\Pwidth}			% for use in tables
\newlength{\Pwidtha}

%
% User-defined environments
%
% "Definition" list with variable-length entries. Requires
% mdwlist package. (See LaTeX Companion, p 64, and mdwlist documentation.)
\newenvironment{Ventry}[1]%
   {\settowidth{\tmplength}{#1\hspace{1em}}%
    \begin{basedescript}{%
           \desclabelwidth{\tmplength}%
           \desclabelstyle{\nextlinelabel}%
           \renewcommand{\makelabel}[1]{##1\hfil}}}
   {\end{basedescript}}
% "Definition" list with variable-length entries, indented. Requires
% mdwlist package. (See LaTeX Companion, p 64, and mdwlist documentation.)
\newenvironment{Ventryi}[1]%
   {\settowidth{\tmplength}{\hspace{\saveparindent}#1\hspace{1em}}%
    \begin{basedescript}{%
           \desclabelwidth{\tmplength}%
           \desclabelstyle{\nextlinelabel}%
           \renewcommand{\makelabel}[1]{\hspace{\saveparindent}##1\hfil}}}
   {\end{basedescript}}
% Indented definition list with variable-length entries, compact.
% (see LaTeX Companion, p 64)
%\newenvironment{Ventryic}[1]%
%   {\begin{list}{}{\renewcommand{\makelabel}[1]{\hspace{\saveparindent}##1\hfil}%
%      \settowidth{\labelwidth}{\hspace{\saveparindent}#1\hspace{1em}}%
%      \setlength{\leftmargin}{\labelwidth+\labelsep}
%      \setlength{\parsep}{0in}%
%      \setlength{\itemsep}{0in}}}%
%   {\end{list}}
% Indented definition list with variable-length entries, compact. Requires
% mdwlist package. (See LaTeX Companion, p 64, and mdwlist documentation.)
\newenvironment{Ventryic}[1]%
   {\settowidth{\tmplength}{\hspace{\saveparindent}#1\hspace{1em}}%
    \begin{basedescript}{%
           \desclabelwidth{\tmplength}%
           \desclabelstyle{\nextlinelabel}%
           \renewcommand{\makelabel}[1]{\hspace{\saveparindent}##1\hfil}
           \setlength{\topsep}{0in}%
           \setlength{\parsep}{0in}%
           \setlength{\itemsep}{0in}}}
   {\end{basedescript}}
% Definition list with variable-length entries, compact. Requires
% mdwlist package. (See LaTeX Companion, p 64, and mdwlist documentation.)
\newenvironment{Ventryc}[1]%
   {\settowidth{\tmplength}{#1\hspace{1em}}%
    \begin{basedescript}{%
           \desclabelwidth{\tmplength}%
           \desclabelstyle{\nextlinelabel}%
           \renewcommand{\makelabel}[1]{##1\hfil}
           \setlength{\topsep}{0in}%
           \setlength{\parsep}{0in}%
           \setlength{\itemsep}{0in}}}
   {\end{basedescript}}

% "changes" environment, to identify changed code by color.  The color is
% given by the argument, with a default of red.  When the environment is
% exited, the color is changed to black.
\newenvironment{changes}[1][red]%
   {\color{#1}}
   {\color{black}}

% Adapted from comp.text.tex post by Keith Reckdahl
% (reckdahl@am-sparc7.Stanford.EDU)
% Indents text from the left margin by current value of the list
% length \leftmargin
\newenvironment{indleft}%
   {\begin{list}{}
      {\setlength{\topsep}{0pt}%
       \setlength{\listparindent}{0pt}%
       \setlength{\itemindent}{0pt}%
       \setlength{\parsep}{\parskip}%
      }%
      \item[]}%
   {\end{list}}
% As above, but surrounding an alltt environment
\newenvironment{indlefttt}%
   {\begin{list}{}
      {\setlength{\topsep}{0pt}%
       \setlength{\listparindent}{0pt}%
       \setlength{\itemindent}{0pt}%
       \setlength{\parsep}{\parskip}%
      }%
      \item[]%
      \begin{alltt}}%
   {\end{alltt}\end{list}}
% As above, but also with a tabbing environment
\newenvironment{tabindlefttt}%
   {\begin{list}{}
      {\setlength{\topsep}{0pt}%
       \setlength{\listparindent}{0pt}%
       \setlength{\itemindent}{0pt}%
       \setlength{\parsep}{\parskip}%
      }%
      \item[]%
      \begin{alltt}\begin{tabbing}}%
   {\end{tabbing}\end{alltt}\end{list}}

%
% User-defined commands
%

% Redefine \paragraph as a display heading
\makeatletter
\renewcommand{\paragraph}{\@startsection
   {paragraph}%			% heading name
   {4}%				% heading level
   {0in}%			% indentation of heading
   {0.5\baselineskip}%		% space before heading
   {0.1\baselineskip}%		% space after heading
   {\normalfont\normalsize}}%	% style
\makeatother

% Degree symbol
\newcommand{\degree}[1][]{\ensuremath{{#1}^\circ}\xspace}

% Degree Rankine symbol
\newcommand{\degR}{\textsuperscript{$\circ$}\textnormal{R}\xspace}

\newcommand{\HDF}{HDF5\xspace}
% --- ADFH should have a ``name'', like MLL
% --- SLL (Storage level library), LLL (Low level library),
% --- PLL (Physical level library), NLL (Node), BLL (Below), FLL (File)
% --- GLL (Generic), KLL (Kernel), CLL (Core), WLL (Wouaw), ZLL (Zoo)
\newcommand{\ADFH}{SLL\xspace}
\newcommand{\SLL}{SLL\xspace}

% ############################################################################
\newcommand{\HDFend}{%
\marginpar[\begin{flushright}%
\vskip -1.5em%
$\bigtriangleup$$\bigtriangleup$$\bigtriangleup$%
\end{flushright}
]{%
\begin{flushright}%
\vskip -1.5em%
$\bigtriangleup$$\bigtriangleup$$\bigtriangleup$%
\end{flushright}
}}
%
\newcommand{\HDFstart}{%
\marginpar[\begin{flushright}%
\vskip 0.5em%
$\bigtriangledown$$\bigtriangledown$$\bigtriangledown$%
\end{flushright}
]{%
\begin{flushleft}%
\vskip 0.5em%
$\bigtriangledown$$\bigtriangledown$$\bigtriangledown$%
\end{flushleft}%
}}
%
\newcommand{\HDFline}{%
\marginpar[\begin{flushright}%
\vskip 1em%
{\LARGE $\triangleright$$\triangleright$$\triangleright$}%
\end{flushright}
]{%
\begin{flushleft}%
\vskip 1em%
{\LARGE $\triangleleft$$\triangleleft$$\triangleleft$}%
\end{flushleft}%
}}
% ############################################################################


% Re-def longtable's caption command to use \captionlabelfont from caption pkg
% (basis lifted from longtable.sty)
\makeatletter
\renewcommand{\LT@makecaption}[3]{%
  \LT@mcol\LT@cols c{\hbox to\z@{\hss\parbox[t]\LTcapwidth{%
    \sbox\@tempboxa{{\captionlabelfont #1{#2: }}#3}%
    \ifdim\wd\@tempboxa>\hsize
      {\captionlabelfont #1{#2: }}#3%
    \else
      \hbox to\hsize{\hfil\box\@tempboxa\hfil}%
    \fi
    \endgraf\vskip\baselineskip}%
  \hss}}}
\makeatother

% Make next page odd, with preceding blank page empty (LaTeX Companion, p 93)
\newcommand{\clearemptydoublepage}{\newpage{\pagestyle{empty}\cleardoublepage}}

% Text subscripts, analogous to \textsuperscript command (from comp.text.tex
% post by rf@cl.cam.ac.uk (Robin Fairbairns))
\makeatletter
\DeclareRobustCommand*\tsub[1]{%
  \@tsub{\selectfont#1}}
\def\@tsub#1{%
  {\m@th\ensuremath{_{\mbox{\fontsize\sf@size\z@#1}}}}}
\makeatother

% Text superscripts, but shorter
\newcommand{\tsup}[1]{\textsuperscript{#1}}

% Shortcuts for specific fonts
\newcommand{\bold}[1]{{\normalfont\textbf{#1}}}  % Bold
\newcommand{\ital}[1]{{\normalfont\textit{#1}}}  % Italic
\newcommand{\key} [1]{{\normalfont\texttt{#1}}}  % Fixed font for keywords
\newcommand{\fort}[1]{{\normalfont\texttt{#1}}}  % Fixed font for Fortran stuff

% Shortcuts for math character sizes
\newcommand{\Dstyle}{\displaystyle}
\newcommand{\Tstyle}{\textstyle}
\newcommand{\Sstyle}{\scriptstyle}

% Use section titles as marks (i.e., in headers) instead of subsection titles
\renewcommand{\sectionmark}[1]{\markright{\thesection\ \ #1}}
\renewcommand{\subsectionmark}[1]{}

% "Better" treatment of headings for Appendices  (LaTeX Companion, pp 29,30)
%\renewcommand{\appendix}{%
%   \newcommand{\app}{%
%      \secdef\Appendix\sAppendix}%
%   \newcommand{\subapp}{%
%      \secdef\subAppendix\subsAppendix}%
%   \newcounter{app}%
%   \newcounter{subapp}[app]%
%   \renewcommand{\theapp}{\Alph{app}}%
%   \renewcommand{\thesubapp}{\theapp.\arabic{subapp}}%
%   \renewcommand{\sectionmark}[1]{\markright{\theapp\ \ ##1}}
% }
%\newcommand{\Appendix}[2][?]{%   Complex form
%   \refstepcounter{app}%
%   \addcontentsline{toc}{section}%
%      {\protect\bfseries{\appendixname~\theapp.\ \ } #1}%
%   {\noindent\Large\bfseries\appendixname\ \theapp.\ \ #2\par}%
%   \sectionmark{#1}\vspace{\baselineskip}}
%\newcommand{\sAppendix}[1]{%   Starred form
%   \refstepcounter{app}%
%   {\noindent\Large\bfseries\appendixname\ \theapp.\ \ #1\par}%
%   \vspace{\baselineskip}}
%
%\newcommand{\subAppendix}[2][?]{%   Complex form
%   \refstepcounter{subapp}%
%   \addcontentsline{toc}{subsection}%
%       {\thesubapp\ \ #1}%
%   \vspace{\baselineskip}
%   {\noindent\large\bfseries\thesubapp\ \ #2\par}%
%   \vspace{\baselineskip}}
%\newcommand{\subsAppendix}[1]{%   Starred form
%   \refstepcounter{subapp}%
%   \vspace{\baselineskip}
%   {\noindent\large\bfseries\thesubapp\ \ #1\par}%
%   \vspace{\baselineskip}}

\begin{document}

\pagenumbering{roman}

%\fancyhead[LE]{\bfseries SIDS-to-HDF File Mapping Manual}
%\fancyhead[RO]{\bfseries \rightmark}
\fancyfoot[LE,RO]{\bfseries \thepage}

%\setlength{\droptitle}{1.0in}
\pretitle{\begin{flushleft}\LARGE%
          \includegraphics[width=2.0in]{../logo/cgns_bw}\\*[0.25in]}
\posttitle{\par\end{flushleft}\vskip 1.0em}
\title{{\bfseries CFD General Notation System\\
HDF5 Implementation}\\*[0.25in]
{\Large Document Version 3.1.2\\
CGNS Version 3.1.3}}
\author{}
\date{}
\maketitle
\thispagestyle{empty}

%\title{{\LARGE \textbf{CGNS}}\\*[0.5in]
%The CFD General Notation System\\
%SIDS-to-HDF File Mapping Manual\footnote
%{This is Version \textcolor{red}{1.3 of this document, released ?? ??? 2002}.
%The current PDF version was derived from the MS Word version of the
%document, created by Boeing.}}
%\author{CGNS Project Group}
%\date{}
%\maketitle
%\thispagestyle{empty}

\clearemptydoublepage
\setlength{\parskip}{0ex}		% remove blank space between paragraphs
\thispagestyle{plain}
\tableofcontents
\listoffigures
%\listoftables
\setlength{\parskip}{\saveparskip}	% put it back

\clearemptydoublepage
\pagenumbering{arabic}

\fancyhead[LE]{\bfseries HDF5 Implementation}
\fancyhead[RO]{\bfseries \rightmark}

\section{Introduction}
\label{s:intro}
\thispagestyle{plain}

Advanced Data Format (ADF) is a library of basic database management
and I/O subroutines that implements a relatively simple hierarchical
database.
ADF is written in ANSI C to enhance portability of the software, and the
design of the database allows for portability of files from platform to
platform.
There is also a Fortran interface.
The files are self-describing (i.e., it is possible to browse files and
determine their contents) and extensible (i.e., many different pieces of
software on different platforms may add or modify information).

The routines allow the user to construct a tree structure with their data.
(See \autoref{f:example} on p.~\pageref*{f:example}.)
This structure is very similar to the directory structures of the UNIX
or DOS operating systems.
ADF also allows links between nodes within the same file or to different
files.
This feature works somewhat like ``soft links'' in the UNIX operating
system.
The major difference between the aforementioned directory structures and
ADF is that the nodes not only contain information about their children
(next lower-level nodes) but may also contain data.

The installation package includes the source code to ADF, the Fortran
interface, sample files, and a simple file browser.

\subsection{History of Project}

ADF was developed as part of the CFD
General Notation System(CGNS) project.
The purpose of the CGNS project is to define the data models for
Navier-Stokes based Computational Fluid Dynamics (CFD) technology,
develop a set of standard interface data structures for those data
models, and develop the software that will allow implementation of those
data structures in existing and future CFD analysis tools.
The CGNS system consists of a collection of conventions, and software
conforming to these conventions, for the storage and retrieval of CFD
data.
Adherence to these conventions is intended to facilitate the exchange of
CFD data between sites, between application codes such as solvers and
grid generators, and across computing platforms.

Once the data models (called the \href{../sids/sids.pdf}{Standard
Interface Data Structures}, or SIDS) were defined, a project was started
to write software that could faithfully reproduce that information on
disk.
ADF is the result of that project.
While ADF was developed specifically for the CFD process, it is quite
general and has no built-in knowledge of CFD; therefore, it should
be applicable to storing any type of data that lends itself to a
hierarchical definition.

\subsection{Other Software Interfaces and ADF}

Two other software packages were investigated as part of this project.
The first database interface investigated was the
Hierarchical Data Format (HDF)
developed at the National Center for
Supercomputing Applications at the University of Illinois.
This database system has a large user base and support and has been
in existence more than 6 years with utilities and graphical routines
written with both a C and a Fortran interface.
The limitation of HDF was that it was not truly hierarchical despite its
name.
Any hierarchy has to be built using naming conventions.
Since the CGNS data models indicated a natural hierarchical structure,
it seemed appropriate to develop database software that worked in that
mode by design.
The ADF design considerations are summarized in
\hyperref[s:design]{Appendix~\ref*{s:design}}.

The second was the
Common File Format
(CFF) developed by McDonnell Douglas Aerospace.
CFF is a second-generation database management system that provides a
unifying file structure for CFD data.
The purpose of the Common File is to insulate the user from the myriad
of different computer types that make up computer systems so that the
user or application programmer may process a file from another machine
without performing explicit conversions.
CFF was written in Fortran; however, it was felt that portability and
extensibility could be enhanced using C.
Much of the experience gained from the McDonnell Douglas group is
incorporated into ADF due to the cooperative efforts of personnel from
the CFD group at McDonnell Douglas Aerospace.

\subsection{Organization of Manual}

The main section of this manual explains the basics of the
hierarchical structure of an ADF database.
In addition, the concept of a node as
the basic building block of the hierarchy is developed in detail.
The remaining sections of this manual are extensive.
They provide a glossary of terms and
conventions, as well as information related to the
ADF version releases,
version control numbering, and
architectures that are supported by ADF.
There are two examples, one in Fortran and one in C, in
\hyperref[s:sampleFortran]{Appendix~\ref*{s:sampleFortran}} and
\hyperref[s:sampleC]{Appendix~\ref*{s:sampleC}} respectively, that
implement the structure illustrated in \autoref{f:example} on
p.~\pageref*{f:example}.
These examples should help familiarize the new user with ADF.
Design considerations, file optimization, and portability issues are
also discussed.
The individual ADF core routines are described in detail in
\hyperref[s:subs]{Appendix~\ref*{s:subs}}.
They are categorized into
database-level routines,
data structure and management subroutines,
data query subroutines,
data I/O subroutines, and some
miscellaneous utility subroutines.
Numerous examples are included to clarify the use of each subroutine.
\hyperref[s:errors]{Appendix~\ref*{s:errors}} provides a summary of ADF
error messages, and \hyperref[s:defaults]{Appendix~\ref*{s:defaults}}
lists default values for various parameters and limits on dimensions of
arrays.

%
% last modif: 2004-12-08
%
% NEW CONCEPT: the ADF node is no more a real node, we say now CGNS node,
% because ADF itself exist at the same time with HDF ! 
%
% ****************************************************************************
% ** WE SHOULD NOT DOCUMENT THE MAPPING ASSUMING THE INTERFACE LIBRARY IS ADFH
% ** ANY HDF TOOL WOULD BE ABLE TO PRODUCE AN HDF5 STRUCTURE UNDERSTOOD AS A
% ** CGNS COMPLIANT STORAGE STRUCTURE 
% ****************************************************************************
%
% I don't know if we have to add some implementation details that would have
% an actual impact on interface. For example, the ways we create datasets 
% and dataspaces, always with default properties. Certainly we should locate
% the places where the properties could/would/should be set.
%
% See ADFH.c comments, such as:
%%   /*
%%    * The ADF documentation allows the dimension values to be changed
%%    * without affecting the data, so long as the data type and number
%%    * of dimensions are the same. With HDF5, we could emulate that by
%%    * using extendable data spaces (with chunking). However this only
%%    * allows the data size to increase, not decrease, and coming up
%%    * with a good value for chunking is difficult. Since changing the
%%    * dimension values without rewiting the data is not a common
%%    * operation, I decided to use fixed sizes, then buffer the data
%%    * in these rare cases.
%%    */
%
% Clearly, this is an implementation choice, but it has an impact on
% interface. The question is : ``Can somebody else implement an interface
% AND an HDF5 file compliant without the emulation described above ?''
%

\section{SLL Mapping to \HDF}
\label{s:summary}
\thispagestyle{plain}

% Cheating: almost replace ADF with HDF ;)

The purpose of the current document is to describe the way in which
CFD data is to be represented in an \HDF data tree. To do this, it is necessary
to first describe the \HDF data structure itself in some
detail. Therefore, a conceptual summary of \HDF is given here in order
to make the current document relatively independent, and to allow
the reader to focus on those aspects of \HDF which are essential to
understanding the file mapping.

Any \HDF file with a conformant mapping is \emph{CGNS/\HDF} compliant. The
mapping has been made using a \emph{per-node} basis. Instead of having a new
mapping dedicated to \HDF, we have made a mapping from the ADF nodes
to a set of \HDF groups. While this is not an optimal mapping, the use of 
such an \HDF node allows us to re-use the ADF API without change.

The \HDF documentation should be used as the authoritative
references to resolve any issues not covered by this summary.

\subsection{General Description of \HDF}

% Here, LOW level is a bit scornfull, we may change this term
The \HDF library provides CGNS with a low level storage system.  
This library is developed and maintained by the National Center
for Supercomputing Applications at the University of Illinois at
Urbana-Champaign.
It offers a
free storage system, with support for a very large number of platforms
used in the scientific world.  This system is in charge of all
physical features such as storage, compression, virtual access to data
and many other services also useful for the CGNS users.

\HDF allows a fine tuning of the physical devices, however, the \SLL
use of \HDF forces all \emph{property lists} to \texttt{DEFAULT}
values.

Refer to the HDF5 web site at \url{http://hdf.ncsa.uiuc.edu/HDF5} for
further information about \HDF.

\subsection{Dedicated \HDF Structures}
\label{s:adfnodestructure}

The \emph{CGNS node} is an \HDF group. It contains attributes, an optional
dataset and optional child groups. This structure is required and can be
extended by other implementors as long as the following requirements are
fullfilled.

In a CGNS tree, most of the nodes are \emph{normal} nodes. These nodes
can optionally contain data and have other nodes as children. 

There are special nodes, such as the root node and
the nodes managing the \emph{links}. These special link structures use
the \HDF mount system. It is not necessary for either the MLL 
or \SLL users to understand the \HDF mount system, but it should be 
understood by \SLL implementors.

\subsubsection{The CGNS Node Mapping}

The CGNS node uses \emph{groups} and \emph{attributes} elements of \HDF. 
A basic CGNS node is composed of elements listed in \autoref{t:groupnode}.
The attribute names have been enclosed in square brackets, because the
SLL uses private attributes with a leading \textit{space} character.
These brackets are not part of the names.

\begin{table}[htbp]
\centering
\caption[Basic Node]{\textbf{Basic Node}}
\label{t:groupnode}
\begin{tabular}{l >{\quad}l >{\quad}l}
\\ \hline\hline \\*[-2ex]
\bold{GROUP} &  &  
\\*[1ex] \hline\hline \\*[-2ex]
% is there a reason why types are not native ints/chars ?
REQUIRED ATTRIBUTE  & \fort{[name]}   & \fort{U8LE[33]} \\
REQUIRED ATTRIBUTE  & \fort{[label]}  & \fort{U8LE[33]} \\
REQUIRED ATTRIBUTE  & \fort{[type]}   & \fort{U8LE[3]} \\
REQUIRED ATTRIBUTE  & \fort{[ order]} & \fort{H5T\_NATIVE\_INT} \\
OPTIONAL DATASET    & \fort{[ data]}  & 
\\*[1ex] \hline\hline
\end{tabular}
\end{table}

All the children groups of a node are understood as \SLL children nodes of this
node, unless the group name starts with a \emph{space} character.
The \emph{dataset} contains the actual data of a node. In the case of
a node without data (i.e. \texttt{MT} type), the dataset does not exist.

In addition to the attributes of a \emph{normal} node, special attribute 
names are reserved for the root node (see 
\autoref{t:rootnode}). 
The name of the root node is \texttt{[HDF5 MotherNode]}, its label is
\texttt{[Root Node of HDF5 File]}, type is \texttt{MT}.

\begin{table}[htbp]
\centering
\caption[Root Node]{\textbf{Root Node}}
\label{t:rootnode}
\begin{tabular}{l >{\quad}l >{\quad}l}
\\ \hline\hline \\*[-2ex]
\bold{GROUP} & \fort{[/]} &  
\\*[1ex] \hline\hline \\*[-2ex]
% is there a reason why types are not native ints/chars ?
REQUIRED ATTRIBUTE  & \fort{[ version]} & \fort{U8LE[33]} \\
REQUIRED ATTRIBUTE  & \fort{[ format]}  & \fort{U8LE[33]}
\\*[1ex] \hline\hline
\end{tabular}
\end{table}

The node hosting the link entries is a special node named \texttt{[ mount]}
without attributes nor data. This group should be a child of the root node,
its absolute path is \texttt{[/ mount]}. The group \texttt{[ mount]} contains
one group for each mounted file refered to by a link. 
Each entry has the attributes described in \autoref{t:linknode}. 
The link node itself is a group with the special name \texttt{[ link]}. 
This implementation is discussed further in \autoref{s:linkmechs}.
% clear enough ?

\begin{table}[htbp]
\centering
\caption[Link Node]{\textbf{Link Node}}
\label{t:linknode}
\begin{tabular}{l >{\quad}l >{\quad}l}
\\ \hline\hline \\*[-2ex]
\bold{GROUP} & \fort{[/ mount/ *]} &  
\\*[1ex] \hline\hline \\*[-2ex]
REQUIRED ATTRIBUTE  & \fort{[ refcnt]} & \fort{H5T\_NATIVE\_INT} \\
REQUIRED ATTRIBUTE  & \fort{[ file]}   & \fort{H5T\_NATIVE\_CHAR}
\\*[1ex] \hline\hline
\end{tabular}
\end{table}

\subsubsection{The Node ID}
The node ID is a unique identifier assigned to each existing node by
\HDF when the file containing it is opened, and to new nodes as they are
created. \ADFH inquiries generally return node IDs as a result and
accept node IDs as input. By building a table of IDs, calling software
can subsequently access specific nodes without further search. The Node
ID is real and is not under user control, its lifetime is ended by
the closing of the \HDF file.

\subsubsection{The Node Name}
The node Name is a 32-byte character field which is user
controllable. Its general use is to distinguish among the children of a
given node; consequently, no two children of the same parent may have
the same Name. Constraints related to the node name are detailed in
\autoref{s:nodename}.

\subsubsection{The Label}
The Label is a 32-byte character field which is user controllable. \ADFH
assigns no formal role to the Label, but the intent was to identify
the structure of the included data. It is common for the various
children of a single parent to store different instances of the same
structure. Therefore, there is no prohibition against more than one
child of the same parent having the same Label.

\subsubsection{The Data Type}
The Data Type is a 32-byte character field which specifies the type and
precision of any data which is stored in the data field. Types provided
by \HDF are listed in \autoref{t:datatype}.

\begin{table}[htbp]
\centering
\caption[Data Types]{\textbf{Data Types}}
\label{t:datatype}
\begin{tabular}{l >{\quad}l >{\quad}l}
\\ \hline\hline \\*[-2ex]
\bold{Data Type} & \bold{Notation} & \bold{\HDF Type}
\\*[1ex] \hline\hline \\*[-2ex]
No Data             & \fort{MT} & \fort{-} \\
Integer 32          & \fort{I4} & \fort{H5T\_NATIVE\_INT32} \\
Integer 64          & \fort{I8} & \fort{H5T\_NATIVE\_INT64} \\
Unsigned Integer 32 & \fort{U4} & \fort{H5T\_NATIVE\_UINT32} \\
Unsigned Integer 64 & \fort{U8} & \fort{H5T\_NATIVE\_UINT64} \\
Real 32             & \fort{R4} & \fort{H5T\_NATIVE\_FLOAT} \\
Real 64             & \fort{R8} & \fort{H5T\_NATIVE\_DOUBLE} \\
Complex 64          & \fort{X4} & \fort{-} \\
Complex 128         & \fort{X8} & \fort{-} \\
Character           & \fort{C1} & \fort{H5T\_NATIVE\_CHAR} \\
Byte                & \fort{B1} & \fort{H5T\_NATIVE\_UCHAR} \\
Link                & \fort{LK} & \fort{-} 
\\*[1ex] \hline\hline
\end{tabular}
\end{table}

\noindent
%If the data type is \fort{MT} or \fort{LK}, the node attributes which
%described the data may be left undefined. 
There is no mapping to \HDF
\fort{MT} and \fort{LK} types, because there is no actual
\emph{data space} associated with the nodes. The type information itself
is stored in the node as the strings ``\fort{MT}'' and ``\fort{LK}''.
The types \fort{X4} and \fort{X8}
are also not mapped.

The data storage format is translated as described in \autoref{t:nativeformat}.
\begin{table}[htbp]
\centering
\caption[Native Formats]{\textbf{Native Formats}}
\label{t:nativeformat}
\begin{tabular}{l >{\quad}l}
\\ \hline\hline \\*[-2ex]
\bold{Native format} & \bold{\HDF Type}
\\*[1ex] \hline\hline \\*[-2ex]
\fort{H5T\_IEEE\_F32BE} & \fort{IEEE\_BIG\_32} \\
\fort{H5T\_IEEE\_F32LE} & \fort{IEEE\_LITTLE\_32} \\
\fort{H5T\_IEEE\_F64BE} & \fort{IEEE\_BIG\_64} \\
\fort{H5T\_IEEE\_F64LE} & \fort{IEEE\_LITTLE\_64}
\\*[1ex] \hline\hline
\end{tabular}
\end{table}

\subsubsection{The Number of Dimensions}

The Data portion of a node is designed to store multi-dimensional
arrays of data, each element of which is presumed to be of the Data
Type specified. The Number of Dimensions specifies the number of
integers required to reference a single datum within the array.

\subsubsection{The Dimension Values}

The Dimension Values are a list of integers expressing the actual
sizes of the stored array in each of the dimensions specified. These
dimensions are stored by the \emph{dataspace} associated
with the \emph{dataset}; no other attributes contains these values.

\subsubsection{The Data}

In an \HDF node, the portion of the node holding the actual stored data array
is a \textit{dataset}.

\subsubsection{The Child Table}

All groups of the current group are said to be the children of this group,
except the groups with a name starting with a space character. Each of
these children groups is a \emph{normal} node (i.e. group).

Children may be identified by their names and labels, and, thence,
by their node IDs once these have been determined. \HDF provides no
notion of order among children, but the \SLL layer adds a creation order
stored in the \texttt{[ order]} attribute. This order is guaranteed to
be the same from call to call, even after the file has been closed and
re-opened.

%In particular, the order of a list of
%children returned by ADF has nothing to do with the order in which they
%were inserted in the file. However, the order returned is consistent
%from call to call provided the file has not been closed and the node
%structure has not been modifed.

Note that there is no \emph{parent} table; that is, a node has no direct
knowledge of its parent. Since calling software must open the file
from the root, it presumably cannot access a child without having
first accessed the parent.  It is the responsibility of the calling
software to record the node ID of the parent if this information will be
required.

\subsubsection{The link mechanism}
\label{s:linkmechs}

A \texttt{LK} typed node is a \emph{link}. Such a node refers to
another node elsewhere. In other words, the \emph{link} has no child
or contents, but is a name of a node somewhere in the current file or
in another file.  The \texttt{LK} typed node is said to be the
\emph{emph} node, while the node elsewhere is said to be the
\emph{destination} node.

Is is the role of the \SLL layer to insure consistency\footnote{In
particular in order to avoid an acyclic graph.} and
transparency of the link mechanism, so that any \emph{normal} node
request to the \emph{link} node is performed as if it is performed
on the destination node.

A \emph{link} destination can be in the same file as the \emph{link} source
or in another file. In both cases, the \emph{link} is made using an \HDF
soft link from source to destination. In the case of a destination in another
file, the destination file is mounted in the \texttt{[/ mount]}\footnote{Absolute path.} group and
the soft link is made from the source to the destination now present in
this mounted file. 
% Ok. Is there still somebody in the room ?
% No ? Now, going to explain how to cook the pancakes.

Each time a file is mounted, the \texttt{[/ refcnt]}\footnote{Absolute path.}
attribute is incremented.
The file is unmounted if there is no reference to itself.



\section{General Remarks}
\label{s:general}
\thispagestyle{plain}

\subsection{Acquiring the Software and Documentation}

The CGNS Mid-Level Library may be downloaded from SourceForge, at
{\itshape\url{http://sourceforge.net/projects/cgns}}.
This manual, as well as the other CGNS documentation, is available in
both HTML and PDF format from the CGNS documentation web site, at
{\itshape\url{http://www.grc.nasa.gov/www/cgns/CGNS_docs_current/}}.

\subsection{Organization of This Manual}

The sections that follow describe the Mid-Level Library functions
in detail.
The first three sections cover some basic file operations (i.e.,
opening and closing a CGNS file, and some configuration options)
(\autoref{s:fileops}), accessing a specific node in a CGNS database
(\autoref{s:navigating}), and error handling (\autoref{s:error}).
The remaining sections describe the functions used to read, write, and
modify nodes and data in a CGNS database.
These sections basically follow the organization used in
the ``Detailed CGNS Node Descriptions'' section of the
\href{../filemap/filemap.pdf}{SIDS File Mapping} manual.

At the start of each sub-section is a \textit{Node} line, listing the
the applicable CGNS node label.

Next is a table illustrating the syntax for the Mid-Level Library
functions.
The C functions are shown in the top half of the table, followed by the
corresponding Fortran routines in the bottom half of the table.
Input variables are shown in an \textcolor{input}{\texttt{upright blue}}
font, and output variables are shown in a
\textcolor{output}{\texttt{\textit{slanted red}}} font.
As of Version 3.1, some of the arguments to the Mid-Level Library have
changed from \texttt{int} to \texttt{cgsize\_t} in order to support 64-bit data.
For each function, the right-hand column lists the modes (read, write,
and/or modify) applicable to that function.

The input and output variables are then listed and defined.

\subsection{Language}

The CGNS Mid-Level Library is written in C, but each function has a
Fortran counterpart.
All function names start with ``\texttt{cg\_}''.
The Fortran functions have the same name as their C counterpart with the
addition of the suffix ``\texttt{\_f}''.

\subsection{Character Strings}

All data structure names and labels in CGNS are limited to 32
characters.
When reading a file, it is advised to pre-allocate the
character string variables to 32 characters in Fortran, and 33
in C (to include the string terminator).
Other character strings, such as the CGNS file name or
descriptor text, are unlimited in length.
The space for unlimited length character strings will be created by the
Mid-Level Library; it is then the responsibility of the application to
release this space by a call to \texttt{cg\_free}, described in
\autoref{s:free}.

\subsection{Error Status}

All C functions return an integer value representing the error
status.
All Fortran functions have an additional parameter, \texttt{ier},
which contains the value of the error status.
An error status different from zero implies that an error occured.
The error message can be printed using the error
handling functions of the CGNS library, described in \autoref{s:error}.
The error codes are coded in the C and Fortran include files
\textit{cgnslib.h} and \textit{cgnslib\_f.h}.

\subsection{Typedefs}
\label{s:typedefs}

Beginning with version 3.1, two new typedef variables have been
introduced to support 64-bit mode. The \texttt{cglong\_t} typedef
is always a 64-bit integer, and \texttt{cgsize\_t} will be either
a 32-bit or 64-bit integer depending on how the library was built.
Many of the C functions in the MLL have been changed
to to use \texttt{cgsize\_t} instead of \texttt{int} in the arguments.
These functions include any that may exceed the \texttt{2Gb} limit
of an \texttt{int}, e.g. zone dimensions, element data, boundary
conditions, and connectivity. In Fortran, all integer data is taken
to be \texttt{integer*4} for 32-bit and \texttt{integer*8} for 64-bit
builds.

Several types of variables are defined using typedefs in the
\textit{cgnslib.h} file.
These are intended to facilitate the implementation of
CGNS in C.
These variable types are defined as an enumeration of key
words admissible for any variable of these types.
The file \textit{cgnslib.h} must be included in any C application
programs which use these data types.

In Fortran, the same key words are defined as integer parameters in the
include file \textit{cgnslib\_f.h}.
Such variables should be declared as \texttt{integer} in Fortran
applications.
The file \textit{cgnslib\_f.h} must be included in any
Fortran application using these key words.

Note that the first two enumerated values in these lists,
\textit{xxxNull} and \textit{xxxUserDefined}, are only available
in the C interface, and are provided in the advent that your
C compiler does strict type checking. In Fortran, these values
are replaced by the numerically equivalent \texttt{CG\_Null}
and \texttt{CG\_UserDefined}. These values are also defined in the
C interface, thus either form may be used. The function prototypes
for the MLL use \texttt{CG\_Null} and \texttt{CG\_UserDefined},
rather than the more specific values.

The list of enumerated values (key words) for each of these variable
types (typedefs) are:
        
{\raggedright
\begin{Ventryi}{\texttt{AverageInterfaceType\_t}}
   \item [\texttt{ZoneType\_t}]
         \texttt{ZoneTypeNull, ZoneTypeUserDefined, Structured, Unstructured}
   \item [\texttt{ElementType\_t}]
         \texttt{ElementTypeNull, ElementTypeUserDefined, NODE, BAR\_2, BAR\_3, TRI\_3, TRI\_6, QUAD\_4,
         QUAD\_8, QUAD\_9, TETRA\_4, TETRA\_10, PYRA\_5,
         PYRA\_14, PENTA\_6, PENTA\_15, PENTA\_18, HEXA\_8,
         HEXA\_20, HEXA\_27, MIXED, PYRA\_13, NGON\_n, NFACE\_n}
   \item [\texttt{DataType\_t}]
         \texttt{DataTypeNull, DataTypeUserDefined, Integer, RealSingle, RealDouble, Character}
   \item [\texttt{DataClass\_t}]
         \texttt{DataClassNull, DataClassUserDefined, Dimensional, NormalizedByDimensional,
         NormalizedByUnknownDimensional, NondimensionalParameter,
         DimensionlessConstant}
   \item [\texttt{MassUnits\_t}]
         \texttt{MassUnitsNull, MassUnitsNullUserDefined, Kilogram, Gram, Slug, PoundMass}
   \item [\texttt{LengthUnits\_t}]
         \texttt{LengthUnitsNull, LengthUnitsUserDefined, Meter, Centimeter, Millimeter, Foot,
         Inch}
   \item [\texttt{TimeUnits\_t}]
         \texttt{TimeUnitsNull, TimeUnitsUserDefined, Second}
   \item [\texttt{TemperatureUnits\_t}]
         \texttt{TemperatureUnitsNull, TemperatureUnitsUserDefined, Kelvin, Celsius, Rankine,
         Fahrenheit}
   \item [\texttt{AngleUnits\_t}]
         \texttt{AngleUnitsNull, AngleUnitsUserDefined, Degree, Radian}
   \item [\texttt{ElectricCurrentUnits\_t}]
         \texttt{ElectricCurrentUnitsNull, ElectricCurrentUnitsUserDefined, Ampere, Abampere, Statampere, Edison,
         auCurrent}
   \item [\texttt{SubstanceAmountUnits\_t}]
         \texttt{SubstanceAmountUnitsNull, SubstanceAmountUnitsUserDefined, Mole, Entities, StandardCubicFoot,
         StandardCubicMeter}
   \item [\texttt{LuminousIntensityUnits\_t}]
         \texttt{LuminousIntensityUnitsNull, LuminousIntensityUnitsUserDefined, Candela, Candle, Carcel, Hefner, Violle}
   \item [\texttt{GoverningEquationsType\_t}]
         \texttt{GoverningEquationsTypeNull, GoverningEquationsTypeUserDefined, FullPotential, Euler, NSLaminar,
         NSTurbulent, NSLaminarIncompressible,
         NSTurbulentIncompressible}
   \item [\texttt{ModelType\_t}]
         \texttt{ModelTypeNull, ModelTypeUserDefined, Ideal, VanderWaals, Constant,
         PowerLaw, SutherlandLaw, ConstantPrandtl, EddyViscosity,
         ReynoldsStress, ReynoldsStressAlgebraic,
         Algebraic\_BaldwinLomax, Algebraic\_CebeciSmith,
         HalfEquation\_JohnsonKing, OneEquation\_BaldwinBarth,
         OneEquation\_SpalartAllmaras, TwoEquation\_JonesLaunder,
         TwoEquation\_MenterSST, TwoEquation\_Wilcox,
	 CaloricallyPerfect, ThermallyPerfect,
	 ConstantDensity, RedlichKwong,
	 Frozen, ThermalEquilib, ThermalNonequilib,
	 ChemicalEquilibCurveFit, ChemicalEquilibMinimization,
	 ChemicalNonequilib, EMElectricField, EMMagneticField,
         EMConductivity, Voltage, Interpolated,
         Equilibrium\_LinRessler, Chemistry\_LinRessler}
   \item [\texttt{GridLocation\_t}]
         \texttt{GridLocationNull, GridLocationUserDefined, Vertex, IFaceCenter, CellCenter, JFaceCenter,
         FaceCenter, KFaceCenter, EdgeCenter}
   \item [\texttt{GridConnectivityType\_t}]
         \texttt{GridConnectivityTypeNull, GridConnectivityTypeUserDefined, Overset, Abutting, Abutting1to1}
   \item [\texttt{PointSetType\_t}]
         \texttt{PointSetTypeNull, PointSetTypeUserDefined, PointList, PointRange, PointListDonor, PointRangeDonor,
         ElementList, ElementRange, CellListDonor}
   \item [\texttt{BCType\_t}]
         \texttt{BCTypeNull, BCTypeUserDefined, BCAxisymmetricWedge,
         BCDegenerateLine, BCExtrapolate,
         BCDegeneratePoint, BCDirichlet, BCFarfield, BCNeumann,
         BCGeneral, BCInflow, BCOutflow, BCInflowSubsonic,
         BCOutflowSubsonic, BCInflowSupersonic,
         BCOutflowSupersonic, BCSymmetryPlane, BCTunnelInflow,
         BCSymmetryPolar, BCTunnelOutflow, BCWallViscous,
         BCWall, BCWallViscousHeatFlux, BCWallInviscid,
         BCWallViscousIsothermal, FamilySpecified}
   \item [\texttt{BCDataType\_t}]
         \texttt{BCDataTypeNull, BCDataTypeUserDefined, Dirichlet, Neumann}
   \item [\texttt{RigidGridMotionType\_t}]
         \texttt{RigidGridMotionTypeNull, RigidGridMotionTypeUserDefined, ConstantRate, VariableRate}
   \item [\texttt{ArbitraryGridMotionType\_t}]
         \texttt{ArbitraryGridMotionTypeNull, ArbitraryGridMotionTypeUserDefined, NonDeformingGrid, DeformingGrid}
   \item [\texttt{SimulationType\_t}]
         \texttt{SimulationTypeNull, SimulationTypeUserDefined, TimeAccurate, NonTimeAccurate}
   \item [\texttt{WallFunctionType\_t}]
         \texttt{WallFunctionTypeNull, WallFunctionTypeUserDefined, Generic}
   \item [\texttt{AreaType\_t}]
         \texttt{AreaTypeNull, AreaTypeUserDefined, BleedArea, CaptureArea}
   \item [\texttt{AverageInterfaceType\_t}]
         \texttt{AverageInterfaceTypeNull, AverageInterfaceTypeUserDefined, AverageAll, AverageCircumferential, AverageRadial, AverageI,
	 AverageJ, AverageK}
\end{Ventryi}
}

\subsection{Character Names For Typedefs}

The CGNS library defines character arrays which map the typedefs above
to character strings.
These are global arrays dimensioned to the size of each list of
typedefs.
To retrieve a character string representation of a typedef, use the
typedef value as an index to the appropiate character array.
For example, to retrieve the string ``\texttt{Meter}''
for the \texttt{LengthUnits\_t Meter} typedef, use
\texttt{LengthUnitsName[Meter]}.
Functions are available to retrieve these names without the need for
direct global data access.
These functions also do bounds checking on the input, and if out of
range, will return the string ``\texttt{<invalid>}''.
An additional benefit is that these will work from within a Windows DLL,
and are thus the recommended access technique.
The routines have the same name as the global data arrays, but with a
``\texttt{cg\_}'' prepended.
For the example above, use ``\texttt{cg\_LengthUnitsName(Meter)}''.

\begin{table}
\begin{center}
\begin{tabular}{l}
\hline
Typedef Name Access Functions \\
\hline
\textcolor{output}{const char *name} = cg\_MassUnitsName(\textcolor{input}{MassUnits\_t type}); \\
\textcolor{output}{const char *name} = cg\_LengthUnitsName(\textcolor{input}{LengthUnits\_t type}); \\
\textcolor{output}{const char *name} = cg\_TimeUnitsName(\textcolor{input}{TimeUnits\_t type}); \\
\textcolor{output}{const char *name} = cg\_TemperatureUnitsName(\textcolor{input}{TemperatureUnits\_t type}); \\
\textcolor{output}{const char *name} = cg\_ElectricCurrentUnitsName(\textcolor{input}{ElectricCurrentUnits\_t type}); \\
\textcolor{output}{const char *name} = cg\_SubstanceAmountUnitsName(\textcolor{input}{SubstanceAmountUnits\_t type}); \\
\textcolor{output}{const char *name} = cg\_LuminousIntensityUnitsName(\textcolor{input}{LuminousIntensityUnits\_t type}); \\
\textcolor{output}{const char *name} = cg\_DataClassName(\textcolor{input}{DataClass\_t type}); \\
\textcolor{output}{const char *name} = cg\_GridLocationName(\textcolor{input}{GridLocation\_t type}); \\
\textcolor{output}{const char *name} = cg\_BCDataTypeName(\textcolor{input}{BCDataType\_t type}); \\
\textcolor{output}{const char *name} = cg\_GridConnectivityTypeName(\textcolor{input}{GridConnectivityType\_t type}); \\
\textcolor{output}{const char *name} = cg\_PointSetTypeName(\textcolor{input}{PointSetType\_t type}); \\
\textcolor{output}{const char *name} = cg\_GoverningEquationsTypeName(\textcolor{input}{GoverningEquationsType\_t type}); \\
\textcolor{output}{const char *name} = cg\_ModelTypeName(\textcolor{input}{ModelType\_t type}); \\
\textcolor{output}{const char *name} = cg\_BCTypeName(\textcolor{input}{BCType\_t type}); \\
\textcolor{output}{const char *name} = cg\_DataTypeName(\textcolor{input}{DataType\_t type}); \\
\textcolor{output}{const char *name} = cg\_ElementTypeName(\textcolor{input}{ElementType\_t type}); \\
\textcolor{output}{const char *name} = cg\_ZoneTypeName(\textcolor{input}{ZoneType\_t type}); \\
\textcolor{output}{const char *name} = cg\_RigidGridMotionTypeName(\textcolor{input}{RigidGridMotionType\_t type}); \\
\textcolor{output}{const char *name} = cg\_ArbitraryGridMotionTypeName(\textcolor{input}{ArbitraryGridMotionType\_t type}); \\
\textcolor{output}{const char *name} = cg\_SimulationTypeName(\textcolor{input}{SimulationType\_t type}); \\
\textcolor{output}{const char *name} = cg\_WallFunctionTypeName(\textcolor{input}{WallFunctionType\_t type}); \\
\textcolor{output}{const char *name} = cg\_AreaTypeName(\textcolor{input}{AreaType\_t type}); \\
\textcolor{output}{const char *name} = cg\_AverageInterfaceTypeName(\textcolor{input}{AverageInterfaceType\_t type}); \\
\hline
\end{tabular}
\end{center}
\end{table}
 
\clearpage

\subsection{64-bit C Portability and Issues}

If you use the \texttt{cgsize\_t} data type in new code,
it will work in both 32 and 64-bit compilation modes.
In order to support CGNS versions prior to 3.1, you may also
want to add something like this to your code:

\noindent \texttt{\#if CGNS\_VERSION < 3100}  \\
\noindent \texttt{\#define cgsize\_t int}  \\
\noindent \texttt{\#endif}

Existing code that uses \texttt{int} will not work with a
CGNS 3.1 library compiled in 64-bit mode. You may want
to add something like this to your code:

\noindent \texttt{\#if CGNS\_VERSION >= 3100 \&\& CG\_BUILD\_64BIT} \\
\noindent \texttt{\#error does not work in 64 bit mode} \\
\noindent \texttt{\#endif}

\noindent or modify your code to use \texttt{cgsize\_t}.

\subsection{64-bit Fortran Portability and Issues}

All integer arguments in the Fortran interface are taken to be
\texttt{integer*4} in 32-bit mode and \texttt{integer*8}
in 64-bit mode. If you have used default or implicit integers in
your Fortran code, it should port to 64-bit mode in most cases by
simply turning on your compiler option that promotes implicit
integers to \texttt{integer*8}. If you have explicitly defined your
integers as \texttt{integer*4}, your code will not work in 64-bit mode.
In that case, you will either need to change them to \texttt{integer}
(recommended for portability) or \texttt{integer*8}.

A new integer parameter has been added to the \textit{cgnslib\_f.h}
header, \texttt{CG\_BUILD\_64BIT}, which will be set
to 1 in 64-bit mode and 0 otherwise. You may use this parameter
to check at run time if the CGNS library has been compiled in 64-bit
mode or not, as in:

\noindent \texttt{if (CG\_BUILD\_64BIT .ne. 0) then} \\
\indent \texttt{print \*,'will not work in 64-bit mode'} \\
\indent \texttt{stop} \\
\noindent \texttt{endif}

If you are using a CGNS library prior to version 3.1, this parameter will
not be defined and you will need to rely on your compiler initializing all
undefined values to 0 (not always the case) for this test to work.

If your compiler supports automatic promotion of integers, and you use
implicit integers, your code should port to 64-bit with the following exception.

If you use an \texttt{Integer} data type in any routine that takes
a data type specification, and an implicit integer for the data, the code
will fail when compiled in 64-bit mode with automatic integer promotion.
An example of this would be:

\noindent \texttt{integer dim} \\
\noindent \texttt{integer data(dim)} \\
\noindent \texttt{call cg\_array\_write\_f('array',Integer,1,dim,data)} \\

This is because the MLL interprets the \texttt{Integer} data type
as \texttt{integer*4} regardless of the compilation mode. The compiler,
however, has automatically promoted \texttt{data} to be
\texttt{integer*8}. What you will need to do to prevent this problem,
is to either explicitly define \texttt{data} as in:

\noindent \texttt{integer dim} \\
\noindent \texttt{integer*4 data(dim)} \\
\noindent \texttt{call cg\_array\_write\_f('array',Integer,1,dim,data)} \\

\noindent or

\noindent \texttt{integer dim} \\
\noindent \texttt{integer*8 data(dim)} \\
\noindent \texttt{call cg\_array\_write\_f('array',LongInteger,1,dim,data)} \\

\noindent or test on \texttt{CG\_BUILD\_64BIT} as in:

\noindent \texttt{integer dim} \\
\noindent \texttt{integer data(dim)} \\
\noindent \texttt{if (CG\_BUILD\_64BIT .eq. 0) then} \\
\indent  \texttt{call cg\_array\_write\_f('array',Integer,1,dim,data)} \\
\noindent \texttt{else} \\
\indent  \texttt{call cg\_array\_write\_f('array',LongInteger,1,dim,data)} \\
\noindent \texttt{endif}

The last 2 options will only work with CGNS Version 3.1, since
\texttt{LongInteger} and \texttt{CG\_BUILD\_64BIT} are
not defined in previous versions.

You may also need to be careful when using integer constants as arguments
in 64-bit mode. If your compiler automatically promotes integer constants
to \texttt{integer*8}, then there is no problem. This is probably
the case if your compiler supports implicit integer promotion. If not,
then the constants will be \texttt{integer*4}, and your code
will not work in 64-bit mode. In that case you will need to do something like:

\noindent \texttt{integer*8 one,dim} \\
\noindent \texttt{integer*4 data(dim)} \\
\noindent \texttt{one = 1} \\
\noindent \texttt{call cg\_array\_write\_f('array',Integer,one,dim,data)}


\end{document}
